\chapter{主丛、相伴丛与丛联络}
\label{ch:16-bundles}
\ChapterPlan{周第6章（本地页码待核对）；Meinrenken主丛讲义及Wendl丛讲义补充}{正文证明初稿；8题详解与3幅原创图，尚未独立数学审阅}

第6章用过渡矩阵拼出向量丛，第10章在局部标架中写出联络。主丛把“选取标架”本身变成一个几何空间：联络在这个空间上是一个全局一形式，局部矩阵只是沿局部截面的拉回。相伴丛则说明，同一个主丛和联络如何同时控制不同表示中的向量、张量及曲率。

本章底流形和总空间仍为有限维、Hausdorff、第二可数、无边界的光滑实流形；结构群为第15章意义下的李群。向量空间可为实或复，复数情形仅指纤维运算复线性，不预设底空间为复流形。主丛采用右作用，表示采用左作用；两种作用在相伴丛商中以逆矩阵配合。证明只依赖此前的拼接、单位分解、ODE、Frobenius和Stokes，不预用分类空间、示性类或一般自由适当作用的商定理。

\Needspace{17\baselineskip}
\section{向量丛同态与拉回丛}
\subsection{覆盖映射的丛同态及普遍性质}
\begin{definition}
设$E\to M,F\to N$为同一标量域上的向量丛，$f:M\to N$光滑。覆盖$f$的丛同态为光滑映射$A:E\to F$，使$\pi_F A=f\pi_E$，且每个$A_x:E_x\to F_{f(x)}$线性。覆盖恒等映射且每条纤维上可逆者称丛同构。
\end{definition}
局部标架下同态写成$(x,v)\mapsto(f(x),a(x)v)$；光滑性等价于矩阵$a$光滑。若$A_x$逐点可逆，逆矩阵的伴随矩阵公式说明逆映射也光滑。因此“纤维上线性可逆”和“总空间光滑”一起给出光滑丛同构，不需额外猜测逆的正则性。

\begin{proposition}\label{prop:bundle-pullback-universal}
集合
\[
 f^*F=\{(x,v)\in M\times F:f(x)=\pi_F(v)\}
\]
有自然的向量丛结构，纤维为$F_{f(x)}$，并有自然丛同态$q:f^*F\to F$。任意覆盖$f$的$A:E\to F$唯一分解为$q\circ\widetilde A$，其中$\widetilde A:E\to f^*F$覆盖$\mathrm{id}_M$。
\end{proposition}
\begin{proof}
若$F|_U\simeq U\times\mathbb K^r$，则对应部分由$(x,v)\mapsto(x,v_U)$识别为$f^{-1}(U)\times\mathbb K^r$。过渡矩阵是原过渡与$f$的复合，满足第6章拼接条件。得到的拓扑就是$M\times F$的子空间拓扑，可在这些局部乘积中直接检验；因而该集合不只是抽象同构于某个拼接丛。这里$f$可以有临界点或为常值映射，投影$F\to N$的局部乘积已足够，不要求$f$为浸入。

分解只能取$\widetilde A(v)=(\pi_E(v),A(v))$；局部式说明其光滑、纤维线性。反向与$q$复合即可，故也得到唯一性。
\end{proof}
若$h:L\to M$，自然同构$h^*(f^*F)\simeq(fh)^*F$把$(y,(h(y),v))$送到$(y,v)$。常值映射$f(x)=p$给出$f^*F=M\times F_p$；恒等映射的拉回自然同构于原丛。原截面$s$给出$f^*s(x)=(x,s(f(x)))$，但$f^*F$的截面一般不全是这种形式。例如常值映射拉回平凡线丛后有任意系数函数，而原截面的拉回在纤维方向只有常值。

\subsection{商丛、分裂及拉回联络}
\begin{proposition}\label{prop:bundle-split}
向量丛$E\to M$具有光滑正定纤维内积（复情形为Hermite内积）。若$S\subset E$是光滑子丛，则$E/S$自然成为向量丛，且$E\simeq S\oplus S^\perp$、$E/S\simeq S^\perp$。任意光滑向量丛短正合列都光滑分裂，但通常没有天然分裂。
\end{proposition}
\begin{proof}
在局部平凡化中取标准内积，用第8章的局部有限单位分解加权相加。正定型的正权和仍正定，复情形同时保留Hermite对称性。将$S$的局部标架作Gram--Schmidt，再在某点补齐到$E$的基并缩小邻域，使所补截面仍线性无关；继续正交化便给出$S^\perp$的光滑标架。因此两者直和为$E$。

把$[v]\in E_x/S_x$送到$v$的正交投影到$S_x^\perp$，良定义且为线性同构。局部标架说明这些同构光滑，给出商丛结构；换用另一内积时，商映射诱导的恒等识别仍在局部光滑，所以商结构不依赖选择。短正合列$0\to E'\xrightarrow{i}E\xrightarrow{q}E''\to0$中，$i(E')$是常秩子丛，$q|_{i(E')^\perp}$逐纤维可逆，光滑逆给出右逆$j:E''\to E$。这个$j$依赖内积，故只证明存在。
\end{proof}

\begin{proposition}\label{prop:bundle-pullback-connection}
若$F\to N$有联络$\nabla$，则$f^*F$有唯一联络$f^*\nabla$，满足对原丛局部截面$s$及$X\in T_xM$有
\[
 (f^*\nabla)_X(f^*s)=\bigl(x,\nabla_{\dd f_xX}s\bigr).
\]
若原联络矩阵为$A$、曲率为$F_A=\dd A+A\wedge A$，则拉回联络和曲率矩阵分别为$f^*A$、$f^*F_A$。
\end{proposition}
\begin{proof}
原局部标架$e$拉回后仍为标架。对任意截面$t=(f^*e)a$，定义
\[
 (f^*\nabla)_Xt=(f^*e)\bigl(Xa+(f^*A)(X)a\bigr).
\]
若$e'=eG$，则拉回标架换为$(f^*e)(G\circ f)$；第10章的换标架公式经拉回仍成立，故局部定义相容。它显然满足联络规则和所给条件；条件决定在拉回标架上的值，再由Leibniz规则决定任意截面，故唯一。外微分与拉回交换、拉回保持楔积，遂给出曲率结论。此论证没有假定任意拉回丛截面都来自一个原截面。
\end{proof}


\Needspace{17\baselineskip}
\section{主丛与相伴向量丛}
\subsection{主丛、局部截面与过渡次序}
\begin{definition}\label{def:principal-bundle}
主$G$丛是光滑满射$\pi:P\to M$及光滑右作用$(p,g)\mapsto pg$，使每个$x$有邻域$U$和右$G$等变微分同胚$\pi^{-1}(U)\simeq U\times G$，覆盖$\mathrm{id}_U$；乘积上的右作用为$(x,h)g=(x,hg)$。主丛同态为覆盖底映射的光滑等变映射。
\end{definition}
局部乘积立即说明$\pi$为下沉，右作用自由，每条纤维上作用传递；但仅给出“自由作用”还不等于已经构造了光滑主丛。例如第15章的无理绕行提醒我们，轨道空间拓扑必须检查。本章采用局部平凡化作为定义，不以未证明的一般商定理替代它。

局部平凡化的单位截面记$s_i:U_i\to P$。每个$p\in P_x$唯一写成$p=s_i(x)h_i$。反过来，任一光滑局部截面$s$给出等变平凡化$\psi_s(x,h)=s(x)h$：在一个已知平凡化中，$s(x)$对应$(x,a(x))$，所以$\psi_s$对应$(x,a(x)h)$，其逆光滑。

\begin{proposition}\label{prop:principal-cocycle}
在重叠处唯一的光滑$g_{ij}:U_i\cap U_j\to G$由$s_j=s_i g_{ij}$定义，满足
\begin{equation}\label{eq:principal-cocycle}
 g_{ii}=e,\quad g_{ji}=g_{ij}^{-1},\quad
 g_{ik}=g_{ij}g_{jk},\qquad h_i=g_{ij}h_j.
\end{equation}
反过来，给定这样的过渡函数，可以唯一拼出带这些局部截面的主丛。主丛平凡当且仅当存在全局光滑截面。
\end{proposition}
\begin{proof}
在$s_i$的平凡化中读取$s_j$的群分量即得$g_{ij}$，故光滑；把$s_k=s_jg_{jk}$代入得到复合次序。反向在不交并$\bigsqcup_iU_i\times G$中规定$(x,h_j)_j\sim(x,g_{ij}h_j)_i$，这些条件保证为等价关系。各块到商的映射是开嵌入，因为任一开集在另一块的原像由光滑可逆过渡给出。不同底点用$M$分开，同一底点上的不同总空间点在一个局部乘积中分开，故商Hausdorff。底空间可数子覆盖和$G$第二可数给出总空间第二可数，光滑过渡给出流形结构。右乘与过渡中的左乘交换，故作用下降且局部等变。固定局部识别的唯一性与第6章同理。

全局截面给出$M\times G\to P$、$(x,h)\mapsto s(x)h$，局部式说明其为微分同胚。平凡丛有单位截面，故反向成立。向量丛的零截面总存在；主丛有全局截面则必须平凡。
\end{proof}
换截面$s_i'=s_i k_i$时，$g_{ij}'=k_i^{-1}g_{ij}k_j$。给定$f:N\to M$，集合$f^*P=\{(y,p):f(y)=\pi(p)\}$按局部乘积同样构成主丛，右作用只乘$p$。若$\Phi:Q\to P$覆盖$f$，则$q\mapsto(\pi_Q(q),\Phi(q))$给出$Q\simeq f^*P$：每条纤维上的等变映射是两个自由传递$G$空间之间的双射，局部表达$(y,h)\mapsto(y,a(y)h)$使逆也光滑。

\subsection{标架主丛}
\begin{proposition}\label{prop:principal-frames}
秩$r$实向量丛$E\to M$的标架丛
\[
 \operatorname{Fr}(E)_x=\operatorname{Iso}(\R^r,E_x)
\]
在右作用$u\cdot g=u\circ g$下成为主$\mathrm{GL}(r,\R)$丛。复丛相应使用$\mathbb C^r$和$\mathrm{GL}(r,\mathbb C)$（视为实李群）。
\end{proposition}
\begin{proof}
局部标架$e_i:\mathbb K^r\to E_x$把任意$u$唯一写成$e_i(x)h_i$，其中$h_i\in\mathrm{GL}(r,\mathbb K)$。若$e_j=e_iG_{ij}$，则$h_i=G_{ij}h_j$，这恰为上一命题的主丛过渡。因此拼接出所列总空间的光滑结构及主作用。
\end{proof}
这里$G_{ij}$是\emph{标架}过渡。第6章使用\emph{纤维坐标}过渡$v_j=g_{ji}^{\mathrm{coord}}v_i$，两者关系为$g_{ji}^{\mathrm{coord}}=G_{ij}^{-1}$。同一个字母指标若不说明作用于基还是分量，就容易把本章的公式写反。

若$E$有正定内积，其正交归一标架组成主$\mathrm O(r)$丛；局部Gram--Schmidt保证平凡化存在。复Hermite丛得到主$\mathrm U(r)$丛；若实丛另有定向，则正向正交标架为主$\mathrm{SO}(r)$丛。这些都是在具体局部乘积中构造，不预用一般结构群约化的存在判据。

\subsection{相伴丛及等变函数}
\begin{theorem}\label{thm:associated-bundle}
设$\rho:G\to\mathrm{GL}(V)$为有限维实或复表示。商集合
\begin{equation}\label{eq:associated-convention}
 P\times_\rho V=(P\times V)/\!\sim,\qquad
 (pg,v)\sim(p,\rho(g)v)
\end{equation}
自然成为底空间$M$上的向量丛，记$[p,v]\mapsto\pi(p)$。它的截面与光滑函数$f:P\to V$中满足$f(pg)=\rho(g)^{-1}f(p)$者一一对应。
\end{theorem}
\begin{proof}
等价关系也可写为右作用$(p,v)g=(pg,\rho(g)^{-1}v)$的轨道关系。对$p=s_i(x)h_i$，定义局部纤维坐标
\[
 [p,v]\longmapsto(x,\rho(h_i)v),\qquad
 (x,a)\longmapsto[s_i(x),a].
\]
换代表$(p,v)$为$(pg,\rho(g)^{-1}v)$不改变$\rho(h_i)v$，所以两式互逆。它们对商拓扑是同胚：商映射局部等同于$(x,h,v)\mapsto(x,\rho(h)v)$，这是开映射，因映射$(x,h,v)\mapsto(x,h,\rho(h)v)$为微分同胚后再作开投影。两图之间系数关系为
\begin{equation}\label{eq:associated-coefficients}
 a_i=\rho(g_{ij})a_j.
\end{equation}
故由向量丛拼接构造得到光滑结构；Hausdorff与第二可数也由同一局部乘积及底空间性质保证，无需先证明整个作用适当。

若$f$等变，$s(x)=[p,f(p)]$不依赖$p\in P_x$。反过来，固定$p$时$V\to(P\times_\rho V)_x$、$v\mapsto[p,v]$是线性同构，令$f(p)$为$s(x)$的唯一系数；换$p$为$pg$便给出所列反变换律。局部式同时证明两方向光滑。
\end{proof}

标架丛恢复原丛的自然同构为
\[
 \operatorname{Fr}(E)\times_{\mathrm{GL}(r)}\mathbb K^r\longrightarrow E,
 \qquad[u,v]\longmapsto u(v).
\]
因为$(ug)(v)=u(gv)$，它与\eqref{eq:associated-convention}完全一致。对偶表示$\rho^*(g)=\rho(g)^{-\mathsf T}$、张量表示和外幂表示分别恢复第6章的对偶、张量及外幂丛。取$V=\mathfrak g,\rho=\operatorname{Ad}$得到伴随丛$\operatorname{Ad}P$，其每条纤维带李括号；自同构保持括号使它不依赖代表。

\begin{figure}[htbp]
\centering\DGDiagram{associated-coordinates}

\caption{相伴丛的同一元素在两个局部截面中表示为$[s_i,a_i]=[s_j,a_j]$。主丛截面右乘$g_{ij}$，系数则乘逆表示；过渡方向与第6章一致。}
\label{fig:associated-coordinates}
\end{figure}

\subsection{水平等变形式的下降}
\begin{definition}
$V$值$k$形式$\alpha$在$P$上称为水平的，若任一输入为竖直向量（属于$\ker\dd\pi$）时值为零；称为$\rho$等变的，若$R_g^*\alpha=\rho(g)^{-1}\alpha$。$k=0$时只要求等变。平凡表示的水平不变形式通常称基本形式。
\end{definition}
\begin{proposition}\label{prop:bundle-tensorial-descent}
水平$\rho$等变$V$值$k$形式与$\Omega^k(M;P\times_\rho V)$自然一一对应。
\end{proposition}
\begin{proof}
对$x\in M$及$v_1,\ldots,v_k\in T_xM$，任选$p\in P_x$和满足$\dd\pi(\widetilde v_j)=v_j$的提升，定义
\[
 \beta_x(v_1,\ldots,v_k)=[p,\alpha_p(\widetilde v_1,\ldots,\widetilde v_k)].
\]
改提升只增加竖直分量，被水平性消去；改$p$为$pg$并将提升右移时，形式值乘$\rho(g)^{-1}$，又被相伴等价关系消去。局部截面的拉回给出其光滑系数。反向借$[p,\cdot]$的逆读取$\beta$，用$\dd\pi$作输入即得水平等变形式。这两项显然互逆。
\end{proof}
非平凡表示中，这不是把$\alpha$下降为底空间上的一个固定$V$值形式；不同纤维的值仍须在相伴丛中比较。这个区别对曲率尤其重要。


\Needspace{17\baselineskip}
\section{联络形式、水平提升与曲率}
\subsection{右作用的基本场与主联络}
对$\xi\in\mathfrak g$定义基本场
\[
 \xi^\#(p)=\left.\frac{\dd}{\dd t}\right|_0p\exp(t\xi).
\]
它的流是$R_{\exp(t\xi)}$。在纤维$P_x\simeq G$中，此场对应\emph{左}不变场，所以
\begin{equation}\label{eq:principal-fundamental}
 [\xi^\#,\eta^\#]=[\xi,\eta]^\#,
 \qquad (R_g)_*\xi^\#=(\operatorname{Ad}_{g^{-1}}\xi)^\#.
\end{equation}
后一式也可微分$p\exp(t\xi)g=pg\exp(t\operatorname{Ad}_{g^{-1}}\xi)$得到。右作用的这里是同态号，与第15章\emph{左}作用的正指数基本场反号不同。局部乘积还说明$\xi\mapsto\xi^\#(p)$是$\mathfrak g\simeq V_pP:=\ker\dd\pi_p$的线性同构。

\begin{definition}\label{def:principal-connection}
主联络为$\mathfrak g$值一形式$\omega\in\Omega^1(P;\mathfrak g)$，满足
\begin{equation}\label{eq:principal-connection-axioms}
 \omega(\xi^\#)=\xi,\qquad R_g^*\omega=\operatorname{Ad}_{g^{-1}}\omega.
\end{equation}
第一条称再现竖直生成元，第二条称右等变。水平空间定义为$H_pP=\ker\omega_p$。
\end{definition}
\begin{proposition}\label{prop:principal-horizontal}
主联络与光滑右$G$不变的补分布$TP=HP\oplus VP$一一对应。任意$X\in T_xM$在给定$p\in P_x$处有唯一水平提升$X^h_p$，并且$\dd R_gX^h_p=X^h_{pg}$。
\end{proposition}
\begin{proof}
再现式使$\omega|_{VP}$为逐点同构，故其核是光滑补空间；等变式给出右不变性。反向将$v\in T_pP$分成水平与竖直两项，竖直项唯一为$\xi^\#(p)$，定义$\omega_p(v)=\xi$。光滑性来自光滑直和投影，\eqref{eq:principal-fundamental}给出等变式。水平提升是$\dd\pi|_{HP}$的逆，其光滑性由局部可逆矩阵得到，唯一性给出最后的等式。
\end{proof}

\subsection{局部势、换截面和存在性}
\begin{theorem}\label{thm:principal-local-form}
若$s:U\to P$为局部截面，定义局部联络势$A=s^*\omega\in\Omega^1(U;\mathfrak g)$。在$\psi_s(x,h)=s(x)h$的坐标中
\begin{equation}\label{eq:principal-local-form}
 \psi_s^*\omega=\operatorname{Ad}_{h^{-1}}A+\theta_h,
 \qquad \theta_h=\dd(L_{h^{-1}})_h\dd h.
\end{equation}
换截面$s'=sg$，其中$g:U\to G$光滑，则
\begin{equation}\label{eq:principal-gauge-local}
 A'=\operatorname{Ad}_{g^{-1}}A+g^*\theta.
\end{equation}
矩阵群中$g^*\theta=g^{-1}\dd g$。反过来，任何满足此换元律的一族$A_i$唯一确定主联络。
\end{theorem}
\begin{proof}
对$T_xU$输入，$\psi_s$的微分是$\dd R_h\dd s$，联络值按右等变性为$\operatorname{Ad}_{h^{-1}}A$。对群方向$\zeta\in T_hG$，写成$\dd L_h\xi$，其像为$\xi^\#$，联络值为$\xi=\theta_h(\zeta)$。线性相加给出公式。将$(x,g(x))$拉回即得换截面式，特别要保留$g$随$x$变化产生的第二项。

反向用\eqref{eq:principal-local-form}在每块定义形式。群乘法的微分给出
\[
 (ab)^*\theta=\operatorname{Ad}_{b^{-1}}a^*\theta+b^*\theta
\]
（这里$a,b$可为同一底空间上的群值函数），因此局部式在重叠处相等；各块再现且等变，遂给出全局联络。公式还说明局部势唯一决定$\omega$。
\end{proof}

\Needspace{14\baselineskip}
\begin{theorem}\label{thm:principal-existence}
每个光滑主丛都有主联络。两联络之差是水平、$\operatorname{Ad}$等变的一形式，因而主联络空间是以$\Omega^1(M;\operatorname{Ad}P)$为模型的仿射空间。
\end{theorem}
\begin{proof}
局部乘积上取$A_i=0$得到局部联络$\omega_i$，选从属于该覆盖的局部有限单位分解$\rho_i$，定义
\[
 \omega=\sum_i(\pi^*\rho_i)\omega_i.
\]
每项在其域外延零，支集从属保证光滑。右作用不改变$\pi^*\rho_i$，所以保留等变性；在$\xi^\#$上取值为$(\sum_i\rho_i)\xi=\xi$。故主联络存在，并不要求$G$紧致。两联络相减后再现项抵消，故差水平且等变；由命题\ref{prop:bundle-tensorial-descent}下降到伴随丛值一形式。反向加任意这样的形式仍满足两个公理。
\end{proof}
若$f:N\to M$，自然同态$q:f^*P\to P$使$q^*\omega$成为拉回主丛的联络：其在基本场上的值和等变性直接由$q$等变得到。不能在任意映射下不加说明地把“水平分布的集合原像”当作补分布；一形式的拉回才直接保留所需公理。

\begin{figure}[htbp]
\centering\DGDiagram{principal-horizontal}

\caption{联络在每点将切向量分成水平和竖直部分。水平部分投影为底空间方向，竖直部分由$\omega$读取为李代数元素；斜线仅表示选取了一个补空间，并无天然正交性。}
\label{fig:principal-horizontal}
\end{figure}

\subsection{水平曲线及其全时间存在}
\begin{theorem}\label{thm:principal-transport}
给定紧区间$[a,b]$上的分段光滑路径$\gamma$及$p_0\in P_{\gamma(a)}$，存在唯一分段光滑水平提升$p(t)$，满足$p(a)=p_0$。移动$\mathcal P_\gamma:p_0\mapsto p(b)$为纤维间的等变微分同胚，保持路径拼接和反向，并对有限维光滑路径族光滑依赖。
\end{theorem}
\begin{proof}
在一块中写$p(t)=s(\gamma(t))h(t)$，令$B(t)=A_{\gamma(t)}(\dot\gamma(t))$。水平条件为
\begin{equation}\label{eq:principal-horizontal-ode}
 h^{-1}\dot h=-\operatorname{Ad}_{h^{-1}}B(t),
 \quad\text{即}\quad \dot h=-\dd(R_h)_eB(t).
\end{equation}
矩阵群即$\dot h=-B(t)h$。ODE局部存在唯一性适用。还需证明在整段上存在，不能只因系数光滑就忽略解逃逸。

对时间$t_0$和从$e$出发的方程，紧时间区间中$B(t)$有紧像。时间依赖的局部ODE定理在$e$附近给出一个统一的小时间长度$\delta>0$，可在每个起始时间使用：可在群的一个固定坐标球上对有界系数及其空间导数取统一界，时间端点用单侧区间处理。从任意$h_0$出发的解是从$e$出发的解再右乘$h_0$，因为右不变方程与该右移相容，故同一个$\delta$也适用。有限次延续便覆盖整段；群非紧并不造成这类方程的有限时爆破。

路径紧像可用有限个平凡化域逐段覆盖，在接头处按过渡函数改坐标，唯一性使这些解拼成同一提升。右平移保持水平，故从$p_0g$出发的提升为$p(t)g$，移动等变。反向恢复初值，说明移动可逆；同一唯一性给出拼接。对有限维光滑路径族，在给定参数附近选统一有限分割和局部坐标，再应用ODE参数光滑依赖逐段复合即可。这里没有建立无限维路径空间的光滑结构。
\end{proof}
由此相伴丛上有移动$[p_0,v]\mapsto[\mathcal P_\gamma(p_0),v]$，等变性保证不依赖代表。给定闭路及$p\in P_x$，唯一$h\in G$满足$\mathcal P_\gamma(p)=ph$；换初始点为$pk$时变成$k^{-1}hk$。因此闭路群元素的具体表示依赖纤维中的基点，纤维等变移动本身不依赖这个选择。

\subsection{曲率的水平性与结构方程}
李代数值形式的括号楔积按无阶乘的shuffle求和定义：若$\alpha$次数为$r$、$\beta$次数为$s$，则
\[
 [\alpha\wedge\beta](v_1,\ldots,v_{r+s})
 =\sum_{\sigma\in\operatorname{Sh}(r,s)}\operatorname{sgn}(\sigma)
 [\alpha(v_{\sigma1},\ldots,v_{\sigma r}),
   \beta(v_{\sigma(r+1)},\ldots,v_{\sigma(r+s)})].
\]
按李代数基展开，用楔积与普通Jacobi式逐分量验证带次数的反对称性、Jacobi式及外微分乘法式。特别地$[\omega\wedge\omega](u,v)=2[\omega(u),\omega(v)]$。下文简写$[\alpha,\beta]=[\alpha\wedge\beta]$。

\Needspace{18\baselineskip}
\begin{theorem}\label{thm:principal-curvature}
主联络的曲率定义为
\begin{equation}\label{eq:principal-curvature}
 \Omega=\dd\omega+\tfrac12[\omega,\omega].
\end{equation}
它水平且$\operatorname{Ad}$等变，故下降为$\Omega^2(M;\operatorname{Ad}P)$的元素。局部曲率为
\[
 F_i=s_i^*\Omega=\dd A_i+\tfrac12[A_i,A_i],\qquad
 F_j=\operatorname{Ad}_{g_{ij}^{-1}}F_i.
\]
在矩阵群中$\tfrac12[A_i,A_i]=A_i\wedge A_i$，正好恢复第10章的曲率矩阵。
\end{theorem}
\begin{proof}
对常$\xi\in\mathfrak g$微分等变公理，得到$\mathcal L_{\xi^\#}\omega=-[\xi,\omega]$。由Cartan公式及$\omega(\xi^\#)=\xi$为常值，
\[
 \iota_{\xi^\#}\dd\omega=-[\xi,\omega],\qquad
 \iota_{\xi^\#}\tfrac12[\omega,\omega]=[\xi,\omega].
\]
两项抵消，因此$\Omega$水平。常$g$的拉回与外微分交换，$\operatorname{Ad}_{g^{-1}}$保持李括号，故曲率等变。沿$s_i$拉回得到局部结构式；下降性或换截面的微分中竖直项被消去给出曲率变换律。
\end{proof}

\Needspace{14\baselineskip}
\begin{proposition}\label{prop:principal-bracket-curvature}
底空间向量场$X,Y$的水平提升满足
\begin{equation}\label{eq:principal-bracket-curvature}
 [X^h,Y^h]-[X,Y]^h
 =-\bigl(\Omega(X^h,Y^h)\bigr)^\#.
\end{equation}
右边按点将李代数值变成该点的竖直向量，不表示一个固定常系数基本场。
\end{proposition}
\begin{proof}
相关场的括号仍相关，所以左边竖直。在水平场上$\Omega(X^h,Y^h)=\dd\omega(X^h,Y^h)=-\omega([X^h,Y^h])$，而$\omega([X,Y]^h)=0$。竖直再现同构即给出结论。这个负号与“水平绕行的首项为负曲率”一致。
\end{proof}

\subsection{相伴联络与协变外微分}
记$\rho_*:\mathfrak g\to\operatorname{End}(V)$为表示的微分。若$f:P\to V$表示一个相伴截面，定义
\[
 (\nabla_Xs)(x)=[p,\dd f_p(X^h_p)].
\]
等变性和水平提升的右移性说明其不依赖$p$，并满足联络的两条Leibniz规则。在局部截面$s_i$中，原截面的系数为$a_i=f\circ s_i$，得到
\begin{equation}\label{eq:associated-connection}
 \nabla a_i=\dd a_i+\rho_*(A_i)a_i,\qquad
 R^E=\rho_*(F_i)\quad\text{（局部矩阵）}.
\end{equation}
第一式可直接由等变函数在竖直方向的导数$\dd f(\xi^\#)=-\rho_*(\xi)f$验证：$\dd s_iX$的水平部分为$\dd s_iX-(A_i(X))^\#$，相减反而产生正的$\rho_*(A_i)$。再对$\dd+\rho_*(A_i)$求曲率并使用$\rho_*[\xi,\eta]=[\rho_*\xi,\rho_*\eta]$，便得第二式。

\begin{theorem}\label{thm:principal-covariant-exterior}
对水平$\rho$等变形式，
\[
 D\alpha=\dd\alpha+\rho_*(\omega)\wedge\alpha
\]
仍水平等变，因而定义相伴丛值形式上的协变外微分，且$D^2\alpha=\rho_*(\Omega)\wedge\alpha$。伴随表示时记$D\alpha=\dd\alpha+[\omega,\alpha]$，曲率满足Bianchi式
\begin{equation}\label{eq:principal-bianchi}
 D\Omega=\dd\Omega+[\omega,\Omega]=0.
\end{equation}
\end{theorem}
\begin{proof}
水平性由Cartan公式检验。等变性微分给出$\mathcal L_{\xi^\#}\alpha=-\rho_*(\xi)\alpha$；因$\iota_{\xi^\#}\alpha=0$，$\iota_{\xi^\#}\dd\alpha$恰为这个负项，而收缩$\rho_*(\omega)\wedge\alpha$给出相反的正项。等变性由表示与伴随作用的相容性及拉回自然性得到。因此局部拉回为$\dd+\rho_*(A_i)\wedge$，逐项展开其平方，一阶项抵消，余下$\rho_*(F_i)\wedge$，证明平方公式。

对Bianchi式，用$\Omega=\dd\omega+[\omega,\omega]/2$以及带次数反对称性得到$\dd\Omega=[\dd\omega,\omega]$。又$[\omega,\dd\omega]=-[\dd\omega,\omega]$；余下$[\omega,[\omega,\omega]]/2$由次数为一的Jacobi恒等式为零。因此$\dd\Omega+[\omega,\Omega]=0$。此处不需要将任意李代数嵌入矩阵代数才能计算。
\end{proof}

\begin{proposition}\label{prop:frame-connection-equivalence}
向量丛$E$上的线性联络与$\operatorname{Fr}(E)$上的主$\mathrm{GL}(r)$联络自然一一对应。保持给定内积的联络恰对应正交标架丛上的主$\mathrm O(r)$联络；复情形相应为$\mathrm U(r)$。
\end{proposition}
\begin{proof}
标准表示的诱导联络以$A_i$为局部矩阵。反向，第10章换标架律为$A_j=G_{ij}^{-1}A_iG_{ij}+G_{ij}^{-1}\dd G_{ij}$，定理\ref{thm:principal-local-form}因此唯一重建主联络；局部互逆即给出全局对应。正交标架中度量相容恰为$A_i^{\mathsf T}+A_i=0$，即取值于$\mathfrak o(r)$；Hermite情形将转置改为共轭转置。同一重建遂可在$\mathrm O(r)$或$\mathrm U(r)$中进行。
\end{proof}

\subsection{平坦性及其局部范围}
\Needspace{15\baselineskip}
\begin{theorem}\label{thm:principal-flat}
主联络曲率为零，当且仅当每点附近存在水平局部截面；等价地，可在局部势全为零的平凡化中描述联络。这些平凡化之间的过渡函数局部常值。
\end{theorem}
\begin{proof}
若$\Omega=0$，命题\ref{prop:principal-bracket-curvature}说明水平分布封闭于括号。第5章Frobenius定理给出通过$p$的局部水平积分片，维数为$\dim M$，且$\dd\pi$在片上可逆。逆函数定理将其缩成一个底空间开集上的截面$s$，于是$s^*\omega=0$。反向若局部势为零，$F_i=0$，水平等变下降性使全局$\Omega=0$。两个零势截面的换元式变成$g_{ij}^*\theta=0$；Maurer--Cartan形式逐点可逆，故$\dd g_{ij}=0$，在每个连通重叠分支上为常值。
\end{proof}
这只给出局部水平截面，不能推出全局水平截面或闭路移动恒等。第10章已有平坦线丛的绕圈例子；下面用主$\mathrm U(1)$丛将同一现象写成规范不变量。


\Needspace{17\baselineskip}
\section{Hopf丛与规范变换}
\subsection{换截面与主丛自同构}
\begin{definition}
主丛的规范变换是覆盖$\mathrm{id}_M$的主丛自同构$\Phi:P\to P$。这些变换组成群。本章约定$\omega^\Phi=\Phi^*\omega$；按映射复合作乘法，这是联络空间上的右作用，因为$(\omega^\Phi)^\Psi=\omega^{\Phi\circ\Psi}$。
\end{definition}
\begin{proposition}\label{prop:principal-gauge}
任意规范变换唯一写成$\Phi(p)=p\,u(p)$，其中光滑$u:P\to G$满足$u(pg)=g^{-1}u(p)g$。局部令$u_i=u\circ s_i$，则
\begin{align}
 u_j&=g_{ij}^{-1}u_i g_{ij},\label{eq:principal-gauge-compatibility}\\
 A_i^\Phi&=\operatorname{Ad}_{u_i^{-1}}A_i+u_i^*\theta,
 & F_i^\Phi&=\operatorname{Ad}_{u_i^{-1}}F_i.\label{eq:principal-gauge-active}
\end{align}
相伴丛上的同构$T_\Phi[p,v]=[\Phi(p),v]$满足
\[
 \nabla^\Phi=T_\Phi^{-1}\nabla T_\Phi.
\]
\end{proposition}
\begin{proof}
在每条自由传递纤维上$u(p)$唯一；局部乘积中读取群坐标保证光滑。等变式$\Phi(pg)=\Phi(p)g$等价于$pg\,u(pg)=pu(p)g$，即所列共轭条件。反向任何这样的$u$定义等变映射，在局部形如$(x,h)\mapsto(x,u_i(x)h)$，逆光滑，故为自同构。代入$s_j=s_i g_{ij}$得到$u_i$的相容式。

$s_i^*\Phi^*\omega=(s_i u_i)^*\omega$，所以局部换元公式直接给出$A_i^\Phi$。曲率由外微分及括号拉回自然性得到$\Omega^\Phi=\Phi^*\Omega$；它水平，故沿变动的$u_i$求拉回时竖直项消去，只剩伴随共轭。相伴丛同构在$s_i$系数中为$a\mapsto\rho(u_i)a$，因此
\[
 \rho(u_i)^{-1}(\dd+\rho_*(A_i))\rho(u_i)
 =\dd+\rho_*(A_i^\Phi),
\]
给出最后一式。
\end{proof}
在\emph{被动}换截面$s_i\mapsto s_i k_i$时，联络不变，只是系数改变；在\emph{主动}规范变换时，固定截面$s_i$，将联络换成$\Phi^*\omega$。两个局部公式相同，但所比较的几何对象不同。这里使用拉回约定；若另取$(\Phi^{-1})^*\omega$作为群作用，矩阵里的$u$便须相应取逆，不能混写。

规范变换保持闭路移动的共轭类。事实上，$\Phi$将$\omega^\Phi$水平提升送为$\omega$水平提升，所以
\[
 \mathcal P_\gamma^{\omega^\Phi}=\Phi^{-1}\mathcal P_\gamma^\omega\Phi.
\]
若$\mathcal P_\gamma^\omega(p)=ph$，则拉回联络在同一$p$处的闭路元素为$u(p)^{-1}h\,u(p)$。这给出了明确可比较的规范不变量，而不是把局部势的每个分量都叫不变量。

\subsection{Hopf丛的底空间与局部截面}
把$\mathbb C^2$看作实四维空间，令
\[
 S^3=\{(z_0,z_1):|z_0|^2+|z_1|^2=1\},\qquad
 (z_0,z_1)\lambda=(z_0\lambda,z_1\lambda),\quad\lambda\in\mathrm U(1).
\]
记$\mathbb{CP}^1$为$\mathbb C^2$中复直线的集合，$[z_0:z_1]$为非零向量张成的直线；投影$\pi:S^3\to\mathbb{CP}^1$取该直线。先把底空间的光滑结构具体构造出来，不借用下一章的复流形理论。

两张卡为$U_0=\{z_0\ne0\}$上的$w=z_1/z_0$，以及$U_1=\{z_1\ne0\}$上的$\zeta=z_0/z_1$，重叠处$\zeta=1/w$。映射
\begin{equation}\label{eq:hopf-sphere}
 [z_0:z_1]\longmapsto
 \frac{(2\operatorname{Re}(\bar z_0z_1),\,2\operatorname{Im}(\bar z_0z_1),\,
 |z_0|^2-|z_1|^2)}{|z_0|^2+|z_1|^2}
\end{equation}
把它与单位$S^2$识别：在$w=x+iy$中变成
\[
 \left(\frac{2x}{1+x^2+y^2},\frac{2y}{1+x^2+y^2},
 \frac{1-x^2-y^2}{1+x^2+y^2}\right),
\]
其逆为$w=(X+iY)/(1+Z)$，另一个极点用$\zeta$卡。这些表达式证明双射及局部光滑逆，并给底空间Hausdorff、第二可数结构。$w$卡的$\dd x\wedge\dd y$正向与$S^2$外法向定向一致。

局部单位向量截面可取
\begin{equation}\label{eq:hopf-sections}
 s_0(w)=\frac{(1,w)}{\sqrt{1+|w|^2}},\qquad
 s_1(\zeta)=\frac{(\zeta,1)}{\sqrt{1+|\zeta|^2}}.
\end{equation}
每条纤维是这些单位向量的全部相位倍数；局部逆直接读取相位，所以$\pi$确为主$\mathrm U(1)$丛。在$w=re^{i\phi}\ne0$处
\begin{equation}\label{eq:hopf-transition}
 s_1=s_0g_{01},\qquad g_{01}=\frac{|w|}{w}=e^{-i\phi}.
\end{equation}
这个负指数由本书$s_1=s_0g_{01}$及$w=z_1/z_0$的次序确定。

\subsection{全局联络、曲率及非平凡性}
\begin{theorem}\label{thm:hopf-connection}
在$S^3$上限制的一形式
\[
 \omega=\bar z_0\dd z_0+\bar z_1\dd z_1
\]
为$i\R$值主联络。局部势及曲率为
\begin{align}
 A_0&=\frac{\bar w\dd w-w\dd\bar w}{2(1+|w|^2)}
     =i\frac{r^2}{1+r^2}\dd\phi,\label{eq:hopf-potentials}\\
 A_1&=\frac{\bar\zeta\dd\zeta-\zeta\dd\bar\zeta}{2(1+|\zeta|^2)},
 & A_1\big|_w&=-i\frac1{1+r^2}\dd\phi,\notag\\
 F&=\dd A_0=\frac{2i}{(1+x^2+y^2)^2}\dd x\wedge\dd y
   =\frac i2\operatorname{vol}_{S^2}.\label{eq:hopf-curvature}
\end{align}
其中第二个$A_1$表达式只在重叠环域中使用；两极处分别使用各自的光滑笛卡儿表达式。全局有
\begin{equation}\label{eq:hopf-flux}
 \frac1{2\pi i}\int_{S^2}F=1.
\end{equation}
所以Hopf主丛没有全局截面，因而不平凡。
\end{theorem}
\begin{proof}
在$S^3$的切向量上，$\omega+\bar\omega=\dd(|z_0|^2+|z_1|^2)=0$，故值为纯虚数。常相位右乘保持$\omega$；$\xi=it$的基本向量为$it(z_0,z_1)$，代入得$\omega(\xi^\#)=it$，满足两个主联络公理。

把\eqref{eq:hopf-sections}代入，求导归一化分母，得到$A_0$的第一式。$w=re^{i\phi}$给第二式；$\zeta=r^{-1}e^{-i\phi}$则给$A_1|_w$。因此$A_1=A_0-i\dd\phi=A_0+g_{01}^{-1}\dd g_{01}$，恰符合换截面律。群交换使曲率为$\dd\omega$，其水平性与不变性已经由一般定理保证，下降为全局$i\R$值二形式$F$。直接对$A_0$求导得
\[
 F=\frac{2ir}{(1+r^2)^2}\dd r\wedge\dd\phi.
\]
单位球面面积元在此坐标为$4\dd x\wedge\dd y/(1+r^2)^2$（由\eqref{eq:hopf-sphere}的两个导数叉积核对），所以$F=i\operatorname{vol}/2$。极点处光滑表达式保证可用$R\to\infty$的圆盘积分逼近全积分，于是
\[
 \int_{S^2}F=2\pi i\int_0^\infty\frac{2r}{(1+r^2)^2}\dd r=2\pi i.
\]
若有全局截面$s$，则$A=s^*\omega$是整个$S^2$上的一形式，且$F=\dd A$。Stokes使$\int_{S^2}F=0$，矛盾。因此主丛不平凡。该论证没有使用$S^3$的基本群或示性类分类。
\end{proof}

\begin{figure}[htbp]
\centering\DGDiagram{hopf-potentials}

\caption{Hopf丛两张卡的局部势。取$r=\tan(\vartheta/2)$，同一纬线的两种势相差$-i\dd\phi$；曲率为单位球面面积形式的$i/2$倍。图中球面为正投影示意，环箭头标明$\phi$增加方向。}
\label{fig:hopf-potentials}
\end{figure}

标准表示$\rho(\lambda)v=\lambda v$的相伴复线丛可识别为
\[
 L=\{(\ell,v):v\in\ell\subset\mathbb C^2\},\qquad
 [z,c]\longmapsto(\mathbb Cz,cz).
\]
它称为重言线丛。这确实良定义，因为$[z\lambda,c]=[z,\lambda c]$。若$L$有无零截面，将其按环境范数归一化便给出Hopf主丛的截面，所以$L$也不平凡。我们只计算\eqref{eq:hopf-flux}的周期；若以后采用酉联络的Chern形式归一化$iF/(2\pi)$，此模型的积分为$-1$。这与$g_{01}=e^{-i\phi}$及重言线丛一致，不能把$1/(2\pi i)$和$i/(2\pi)$当作同一个系数。一般Chern类理论留到第18章。

\subsection{曲率、闭路相位与平坦反例}
在$w$卡沿$r=r_0$、$0\le\phi\le2\pi$走一圈，写水平提升为$p(\phi)=s_0(r_0e^{i\phi})e^{i\chi(\phi)}$。水平式给
\[
 \chi'=-\frac{r_0^2}{1+r_0^2},\qquad
 h=\exp\left(-2\pi i\frac{r_0^2}{1+r_0^2}\right).
\]
以北极为中心的球帽面积为$\mathcal A=4\pi r_0^2/(1+r_0^2)$，故$h=e^{-i\mathcal A/2}$。改用另一块跨过南极计算时，相位差是$2\pi$的整数倍，最终群元素相同。此处“面积决定相位”是这个明确联络和可张成球帽的闭路的计算，不是对任意主丛联络的普遍公式。

另一方面，在平凡主丛$S^1\times\mathrm U(1)$上，令
\[
 \omega=h^{-1}\dd h+i\kappa\,\dd\theta,\qquad\kappa\in\R.
\]
这里$\dd\theta$指圆周上全局一形式$x\dd y-y\dd x$，不要求有全局实值角函数。曲率恒零，但绕圈移动为$e^{-2\pi i\kappa}$，$\kappa\notin\mathbb Z$时非平凡。与Hopf丛相反，这个\emph{丛}有全局截面，却没有全局\emph{水平}截面；主丛平凡、联络平坦、闭路移动恒等是三个需要分别检验的条件。


\section{主丛与联络习题及详细解答}
16.1改述Wendl第2章原题2.34；16.2改述原题2.90；16.3取原题2.13的光滑版本并扩展复线丛及Hopf例子。原位置分别为印刷33/PDF17、印刷66/PDF50、印刷24/PDF8（均为该章独立PDF）。其余五题为原创补充。周书本地页码和梁书对应题仍待核实，不据主题相近就填入原题号。

\Needspace{12\baselineskip}
\subsection*{习题16.1：对角拉回与切映射的目标丛}
改述Wendl原题2.34并扩展。对$\Delta:M\to M\times M$、$\Delta(x)=(x,x)$，证明$\Delta^*T(M\times M)\simeq TM\oplus TM$。在此识别下写出$\dd\Delta$。一般$f:M\to N$的微分应视为哪两个底空间相同的丛之间的同态？
\begin{proof}[解答]
两因子投影的微分给出自然线性同构$T_{(x,y)}(M\times M)\to T_xM\oplus T_yM$。在乘积坐标中它就是把速度的两组分量分开，逆是把两组速度合并，所以光滑。限制到$y=x$，便得到
\[
 (x,v)\longmapsto(x,(\dd\operatorname{pr}_1v,\dd\operatorname{pr}_2v))
\]
的丛同构。因$\operatorname{pr}_j\Delta=\mathrm{id}$，$\dd\Delta$在此识别下为$X\mapsto(X,X)$，不是任一因子的包含。它的像是直和丛中的对角子丛，商可由$(u,v)\mapsto v-u$识别为$TM$。

对一般$f$，$\dd f:TM\to TN$覆盖$f$，因而唯一因子化为$TM\to f^*TN$、$X_x\mapsto(x,\dd f_xX_x)$，此时两丛底空间都是$M$。不要求$f$为浸入；若微分秩跳跃，其像不能直接称光滑子丛。取$f(t)=t^2$时，原点的像秩零、其余点秩一，正说明这个区别。
\end{proof}

\Needspace{12\baselineskip}
\subsection*{习题16.2：主丛等变映射为何自动可逆}
改述Wendl原题2.90。设$P,Q$为同一底空间上的主$G$丛，$\Phi:P\to Q$光滑、覆盖恒等且等变。证明它自动为主丛同构，写出局部逆。若底映射不是恒等，这个结论如何改写？
\begin{proof}[解答]
固定局部截面$s$和$t$，唯一光滑$a(x)\in G$使$\Phi(s(x))=t(x)a(x)$。因此
\[
 \Phi(s(x)h)=t(x)a(x)h,
 \qquad t(x)k\longmapsto s(x)a(x)^{-1}k
\]
为互逆局部表达。纤维自由传递使这些局部逆不依赖所选截面，故拼成全局光滑逆，并等变。

若$\Phi$覆盖$f:M\to N$，正确结论是$P\simeq f^*Q$，同构为$p\mapsto(\pi(p),\Phi(p))$。不能直接断言$P\simeq Q$或$\Phi$双射。例如$P=Q=S^1\times G$、$\Phi(z,h)=(z^2,h)$为等变主丛同态，但总空间上不是单射；它仍给出$P$与相应拉回丛的同构。
\end{proof}

\Needspace{12\baselineskip}
\subsection*{习题16.3：无零截面与相伴重言线丛}
取Wendl原题2.13的光滑版本：证明实或复光滑线丛平凡当且仅当有光滑无零截面。将它用于Hopf丛的标准表示相伴线丛，说明该线丛非平凡。为什么原题的“连续”版本不能在此被无说明地升级为光滑结论？
\begin{proof}[解答]
若$s$无零，则$(x,c)\mapsto c\,s(x)$给出纤维线性双射。局部标架$e$中写$s=e a(x)$，其中$a(x)\ne0$，逆用$1/a(x)$，故为光滑丛同构。反过来，平凡丛的常截面1经同构移回即得无零截面。实数和复数的证明相同。

Hopf相伴丛的元素$[z,c]$映为直线$\ell=\mathbb Cz$中的向量$cz$。若有无零截面$v(\ell)$，其归一化$v(\ell)/\|v(\ell)\|$为$S^3\to\mathbb{CP}^1$的光滑右逆，这与\eqref{eq:hopf-flux}及Stokes矛盾。所以重言线丛无光滑无零截面，不平凡。

原题对连续丛的连续无零截面给出拓扑平凡化；本题把假设和结论同时改为光滑。若仅给连续截面而要求光滑平凡化，还须证明适当的光滑逼近并保持无零性，不能从连续函数的倒数存在就推断其光滑。本书此处没有使用这项逼近定理。
\end{proof}

\Needspace{12\baselineskip}
\subsection*{习题16.4：Möbius拉回和离散群的平坦移动}
令$L=(\R\times\R)/((t,a)\sim(t+2\pi,-a))$，以$[t,a]\mapsto e^{it}$投影。证明$z\mapsto z^2$拉回$L$后平凡。再取$L$的单位标架主$\mathrm O(1)$丛，求其曲率及绕一圈的闭路移动。
\begin{proof}[解答]
拉回丛中的截面可写为$e^{it}\mapsto(e^{it},[2t,1])$。当$t$换为$t+2\pi$时，$[2t+4\pi,1]=[2t,1]$，所以它良定义、光滑且无零，给出平凡化。更一般地，$m$重映射的拉回为$(t,a)\sim(t+2\pi,(-1)^ma)$；$m$偶数时平凡，$m$奇数时用中间值定理排除无零实截面。

度量$|[t,a]|=|a|$良定义。单位标架$u_t(c)=[t,c]$满足$u_{t+2\pi}=u_t(-1)$，所以标架丛为非平凡二重覆盖。结构群$\mathrm O(1)=\{1,-1\}$的李代数为零，竖直切空间也为零，唯一主联络为零形式，水平空间等于全部切空间，曲率零。绕底圆周的提升从$u_0$走到$u_{2\pi}=-u_0$，闭路元素为$-1$。这是群离散时仍可有非平凡闭路移动的例子；曲率的李代数值信息不能检测这个离散元素。
\end{proof}

\Needspace{12\baselineskip}
\subsection*{习题16.5：非交换规范变换的逐项计算}
在平凡主$\mathrm{GL}(2,\R)$丛上取
\[
 H=\begin{pmatrix}1&0\\0&-1\end{pmatrix},\quad
 E=\begin{pmatrix}0&1\\0&0\end{pmatrix},\quad
 A=xH\dd y,\quad g=I+xE.
\]
计算$A'=g^{-1}Ag+g^{-1}\dd g$及$F'$，逐项核对$F'=g^{-1}Fg$，说明为何不能省略$A'\wedge A'$。
\begin{proof}[解答]
$E^2=0$，所以$g^{-1}=I-xE$；$[H,E]=2E$给$g^{-1}Hg=H+2xE$。因此
\[
 A'=E\dd x+(xH+2x^2E)\dd y.
\]
外微分与乘积项分别为
\[
 \dd A'=(H+4xE)\dd x\wedge\dd y,\qquad
 A'\wedge A'=[E,xH+2x^2E]\dd x\wedge\dd y
 =-2xE\dd x\wedge\dd y.
\]
相加得$F'=(H+2xE)\dd x\wedge\dd y$。原来$F=H\dd x\wedge\dd y$，共轭后恰是同一结果。若只算$\dd A'$，会留下多余的$2xE$，违反曲率变换律。对系数列$a$，也有$(\dd+A')(g^{-1}a)=g^{-1}(\dd+A)a$；这里同时改变截面系数时导数项必参与，不能只把联络矩阵作逐点相似变换。
\end{proof}

\Needspace{12\baselineskip}
\subsection*{习题16.6：Hopf纬线的两种局部势}
沿$w=re^{i\phi}$、$0\le\phi\le2\pi$走一圈，$r>0$固定。分别用$A_0$和$A_1$求水平提升的相位增量，说明为何两个增量不同但闭路群元素相同，并用球帽面积重写结果。
\begin{proof}[解答]
在$s_0$坐标中，水平条件给$\Delta\chi_0=-2\pi r^2/(1+r^2)$。在$s_1$坐标中，$A_1|_w=-i\dd\phi/(1+r^2)$，所以$\Delta\chi_1=2\pi/(1+r^2)$。同一点的坐标关系为$s_0h_0=s_1h_1=s_0e^{-i\phi}h_1$，即$h_1=e^{i\phi}h_0$，两个相位增量相差$2\pi$，因此指数相同。

由球面面积元，北球帽面积$\mathcal A=4\pi r^2/(1+r^2)$，故$h=e^{-i\mathcal A/2}$。用互补球帽并反向取边界，得$e^{i(4\pi-\mathcal A)/2}$，与前式相同。特别地$r=1$时闭路元素为$-1$；纬线缩到北极$r\to0$或南极$r\to\infty$时均趋于1。相位的一个实数提升可随卡改变整周，而圆群元素是内禀的。
\end{proof}

\Needspace{12\baselineskip}
\subsection*{习题16.7：圆群曲率周期及表示的整数权}
设任意主$\mathrm U(1)$丛在定向$S^2$上带联络，曲率下降为$i\R$值二形式$F$。证明其曲率周期满足
\[
 \frac1{2\pi i}\int_{S^2}F\in\mathbb Z.
\]
只用两个圆盘上的截面及绕数的基本计算，不引用线丛分类。再考虑Hopf丛的整数权表示
\[
 \rho_k(\lambda)c=\lambda^kc,\qquad k\in\mathbb Z,
\]
求其相伴线丛的曲率，并判断非平凡性。
\begin{proof}[解答]
先说明两个略超过半球的开圆盘上确有截面。每个圆盘取中心，沿坐标径向路径用主联络作水平提升；定理\ref{thm:principal-transport}给出光滑依赖和全时间存在，故得到光滑截面。分别记$s_N,s_S$，赤道上$s_S=s_Ng$。以北半球边界定向参数化赤道，Stokes及局部换元律给
\[
 \int_{S^2}F=\int_{S^1}(A_N-A_S)=-\int_{S^1}g^{-1}\dd g.
\]
为计算最后一项，把$g(e^{i\phi})$在区间$[0,2\pi]$上写成$e^{i\psi(\phi)}$：先选初始相位，再解$\psi'=-ig^{-1}g'$，求导说明$e^{-i\psi}g$恒为1，因此该写法存在。端点相同迫使$\psi(2\pi)-\psi(0)=2\pi m$，$m\in\mathbb Z$。于是积分周期为$-m$。这只证明周期整数性，没有证明周期给出所有丛的完整分类。

Hopf丛的$g_{NS}=e^{-i\phi}$有$m=-1$，与本章周期$+1$相符。权$k$表示的微分把$it$送为$ikt$，所以相伴线丛联络势为$kA_i$、曲率为$kF$，周期为$k$。$k=0$时表示平凡，相伴丛平凡且诱导联络为平凡联络。若$k\ne0$且线丛平凡，取全局无零截面给出全局联络一形式，其曲率为外微分，积分必须为零，与$2\pi ik$矛盾。因此这些非零权线丛都非平凡；此结论不依赖一般Chern类构造。
\end{proof}

\Needspace{12\baselineskip}
\subsection*{习题16.8：圆周上的平坦联络及大规范变换}
对平凡主$\mathrm U(1)$丛取$A=i\kappa\dd\theta$。证明任一光滑规范变换$u:S^1\to\mathrm U(1)$使势的周期只改变$2\pi i$的整数倍，并判定两个常数势$i\kappa\dd\theta,i\kappa'\dd\theta$规范等价的条件。何时有全局水平截面？
\begin{proof}[解答]
圆群交换，$A^u=A+u^{-1}\dd u$。如上一题，在参数区间上取$u=e^{i\psi}$，端点差为$2\pi m$，所以
\[
 \int_{S^1}(A^u-A)=2\pi im.
\]
若两个常数势规范等价，则$\kappa'-\kappa=m\in\mathbb Z$。反过来，$u(e^{i\theta})=e^{im\theta}$给出$A^u=i(\kappa+m)\dd\theta$，所以该整数条件也充分。

全局截面可写为$e^{i\theta}\mapsto(e^{i\theta},h(\theta))$。水平条件为$h'=-i\kappa h$，解为$h(\theta)=e^{-i\kappa\theta}h(0)$；它在圆周上单值当且仅当$\kappa\in\mathbb Z$。这又等价于闭路移动$e^{-2\pi i\kappa}=1$，也等价于可规范化为零势。所有这些联络的曲率都为零，故仅凭曲率不能区分其规范等价类；这里的常数族由$\R/\mathbb Z$参数化。
\end{proof}

